% Image credits for the photographs and AI illustrations of Book 4
% (University Chemistry, Year 3), Hindi edition. Machine-readable license
% records live in images/book4/CREDITS.md; keep this file in sync with the
% English frontmatter/image-credits-book4.tex. Authors' names stay in their
% original Latin spelling here (attribution), as in the Book 2 Hindi credits.
\chapter*{चित्र-श्रेय}

इस पुस्तक के रेखाचित्र मौलिक हैं। जहाँ छायाचित्र प्रयुक्त हैं, वे अपने-अपने मुक्त
अनुज्ञापत्रों के अंतर्गत प्रयुक्त हैं, उनके रचयिताओं के प्रति आभार सहित; हर छायाचित्र के
स्रोत की पूरी कड़ियाँ और अनुज्ञापत्रों के पते पुस्तक के रिपॉज़िटरी में
\texttt{images/book4/CREDITS.md} में सूचीबद्ध हैं।
संकट चित्रलेख रसायनों के वर्गीकरण और लेबलन की विश्व स्तर पर सुसंगत प्रणाली (GHS) के
चित्रलेख हैं, जैसे उन्हें \TeX\ Live का \texttt{ghsystem} पैकेज खींचता है।

\bigskip
छायाचित्रों और मौलिक रेखाचित्रों के साथ-साथ, कुछ चित्र
कृत्रिम बुद्धि से बनाए गए चित्रण हैं, जो पुस्तक के संपादकों द्वारा लिखे गए निर्देशों से
OpenAI के चित्र-निर्माण द्वारा बनाए गए हैं, और समावेश से पहले रासायनिक
यथार्थता के लिए जाँचे गए हैं। हर निर्मित चित्र का पूरा निर्देश
रिपॉज़िटरी में \texttt{images/book4/ai/PROMPTS.md} में अभिलिखित है।
\bigskip
\noindent छायाचित्र, अध्याय के अनुसार:
\begin{itemize}\small
  \item अध्याय~1: एर्विन श्रोडिंगर, 1933, Nobel Foundation, सार्वजनिक अधिकार-क्षेत्र
        (Wikimedia Commons)।
  \item अध्याय~4: हिम-क्रिस्टल, Wilson A. Bentley (लगभग 1902), सार्वजनिक अधिकार-क्षेत्र
        (Wikimedia Commons)।
  \item अध्याय~6: एक ऊँचे पठार पर रेडियो एंटेना, ESO/B.~Tafreshi
        (twanight.org), CC BY 4.0 (Wikimedia Commons)।
  \item अध्याय~7: दिन के प्रकाश में और पराबैंगनी प्रकाश में फ़्लुओराइट, Masha
        Milshina द्वारा, CC BY 4.0 (Wikimedia Commons)।
  \item अध्याय~8: एक NMR प्रयोगशाला, Shandchem द्वारा, CC BY 2.0 (Wikimedia
        Commons)।
  \item अध्याय~9: विलियम लॉरेंस ब्रैग, Nobel Foundation, सार्वजनिक अधिकार-क्षेत्र
        (Wikimedia Commons)।
  \item अध्याय~10: लुडविग बोल्ट्ज़मान की समाधि, Feldkurat Katz द्वारा,
        CC BY-SA 4.0 (Wikimedia Commons)।
  \item अध्याय~13: बेलूसोव--झाबोटिंस्की सर्पिल, Stephen Morris द्वारा, CC BY 2.0
        (Wikimedia Commons)।
  \item अध्याय~14: अंटार्कटिक ओज़ोन छिद्र, NASA, सार्वजनिक अधिकार-क्षेत्र (Wikimedia
        Commons)।
  \item अध्याय~16: उत्प्रेरकी परिवर्तक के मधुकोश का इलेक्ट्रॉन सूक्ष्मचित्र,
        Galadrid द्वारा, CC BY-SA 4.0 (Wikimedia Commons)।
  \item अध्याय~17: सूर्य-किरणें और टिंडल प्रभाव, Kevin c higgins द्वारा, CC BY-SA 4.0
        (Wikimedia Commons)।
  \item अध्याय~18: माणिक्य और पन्ने के क्रिस्टल, Robert M. Lavinsky द्वारा, CC BY-SA 3.0
        (Wikimedia Commons)।
  \item अध्याय~20: फ़ेरोसीन के क्रिस्टल, Andrea Sella द्वारा, CC BY 4.0 (Wikimedia
        Commons)।
  \item अध्याय~22: एकल-क्रिस्टल सिलिकॉन बूल, Sebastian Wallroth द्वारा, सार्वजनिक
        अधिकार-क्षेत्र (Wikimedia Commons)।
  \item अध्याय~23: सिलिका ऐरोजेल पर एक फूल, NASA/JPL, सार्वजनिक अधिकार-क्षेत्र;
        लाइकर्गस प्याला, Marie-Lan Nguyen द्वारा, CC BY 2.5 (Wikimedia Commons)।
  \item अध्याय~24: एक अटलांटिक हॉर्सशू केकड़ा, Kaldari द्वारा, CC0 (Wikimedia
        Commons)।
  \item अध्याय~27: पाइरेथ्रम डेज़ी, Soramimi द्वारा, CC BY-SA 4.0 (Wikimedia
        Commons)।
  \item अध्याय~30: पैसिफ़िक यू की छाल, JOE BLOWE (theforestprimeval) द्वारा,
        CC BY-SA 2.0 (Wikimedia Commons)।
  \item अध्याय~32: लैंडसैट 8 उपग्रह से देखा गया एक शैवाल-प्रस्फुटन, USGS और NASA,
        सार्वजनिक अधिकार-क्षेत्र।
  \item अध्याय~33: एक दोहरा निर्वात और अक्रिय-गैस नलिका-तंत्र, Polimerek द्वारा, CC BY-SA
        3.0 (Wikimedia Commons)।
\end{itemize}
\clearpage
